\FloatBarrier
\section{Discussion}
\label{sec:discussion}

\paragraph{A recipe for model-assisted coding.} For a researcher who wants to code a corpus with a fixed codebook and check
only part of it by hand, the results suggest four steps.
\begin{enumerate}
  \item Hand-code a random sample of units from the corpus itself, and set part of it aside for certification. Units can be
  comments, transcripts, random fixed-size blocks of turns, or single utterances coded with their context; units of similar
  weight make a guarantee far cheaper (\S\ref{sec:eq2}). Random selection is what makes the guarantee possible; thresholds
  from other codebooks, and fixed rules such as ``accept $\ge$0.95'', did not carry over in our data.
  \item On the other part, a pilot sample, choose the coder and its scores. Compare how many decisions each candidate
  accepts at the target, which after recalibration depends mostly on accuracy; low ECE is not a sufficient criterion
  (\S\ref{sec:eq1}). Fit any per-code maps here: they improved error ranking for every coder with enough calibration data,
  but maps fitted on the certification sample would void its guarantee (\S\ref{sec:eq3}). Compare the accepted share with a
  prevalence baseline that reads no text, and check how many applying codes the accepted decisions contain (\S\ref{sec:rq1}).
  \item Freeze the chosen coder, scores and threshold grid, then certify on the set-aside sample with a rule that respects how
  it was drawn, such as corpus LTT. Each certificate covers one configuration: choosing among configurations by their
  certification results needs all of them in one testing family. Our version needed about 120 error-free units of similar
  weight before it certified anything; with whole transcripts of very unequal length it could not certify even a perfect
  coder (\S\ref{sec:eq2}).
  \item Count what a person still has to read, including the hand-coded sample: in multi-label coding, a segment leaves the
  queue only when all its codes are accepted (\S\ref{sec:eq5}).
\end{enumerate}
In our data, pplx-decider alone gave the most certified automation per inference dollar on comments, and Clef, the pplx-decider +
Kev-27B gate and the decision-model committee certified within two points of Sonnet, at 5--21\% of its inference cost; the
intervals do not resolve those differences. Once human reading is priced in, the ranking can change (\S\ref{sec:eq5}).

\paragraph{What decision models are, and are not.} Decision models are marketed as calibrated. In our data four of the seven
have raw ECEs of 0.05--0.07 and Jev, OpenAI Decisions and Kev-9B 0.10--0.12; a Platt map per task reduces all of them to
0.03--0.07. What separates them raw is error ranking and resolution at the top of the scale, where the category ranges from
worse than a cheap LLM's verbalized confidence (Kev-9B, Clef-flash) to near-frontier (Kev-27B). Per-code recalibration narrows
most of the ranking differences, partly by using which code was predicted, and leaves accuracy as the main differentiator.
The comparison with a 9B instruct model, together with \citet{li2026llmasjev}, is consistent with the hypothesis that much
of a decision model's advantage comes from reading out probabilities rather than generating them. Option order still
changes up to 31\% of a decision model's labels on average over tasks, and 55\% in one task (\S\ref{sec:rq3}); averaging over option rotations \citep{zhang2026anyjev} is
a cheap remedy we did not test.

\paragraph{Human labels are not gold.} No coder fell below the range of individual annotators, and adjudication often sided
with the model on sampled disagreements. Error measured against a single majority label may therefore overstate some model
errors, and a 5\% target against such a reference is partly a target against annotator noise
\citep{plank2022problem,uma2021learning}. Agreement with annotators, which we measure, is not the same as interpretive
validity, which we do not.

\paragraph{From labels to estimates.} Triage bounds the error rate among the labels a researcher accepts, in aggregate; it
says nothing about any one label. Many
qualitative studies go on to report code frequencies or associations between codes. Prediction-powered and design-based
estimators correct such quantities using a hand-coded subsample, without per-item guarantees
\citep{angelopoulos2023ppi,egami2023using}, and model confidence can decide which items to hand-code
\citep{gligoric2025unconfident}. The calibration sample that triage needs can serve these estimators too, and they do not
need a certified threshold, which makes them attractive on dialogue data, where we could certify almost nothing under either
sampling design.

\section{Limitations and follow-ups}
\label{sec:limitations}
\begin{itemize}
  \item \textbf{Scope.} Five tasks on short segments (comments, utterances, sentences) from three public datasets. We have no
  long-form interviews: public interview data with applied codes is rare, and a private dataset requires ethics approval.
  Pooled estimates describe these five tasks, not a population of codebooks.
  \item \textbf{A dated snapshot.} Several APIs were in public beta, and the models will change. We release the harness, the
  analysis plan and every model output, so every number can be recomputed from the public datasets and the study re-run.
  \item \textbf{Guarantees.} Corpus LTT assumes the hand-coded clusters are a uniformly random subset of the corpus and that
  the test visits them in random orders; our split assigns clusters by a hash of their id, which we treat as random, and the
  orders come from a recorded seed. Each certificate covers one configuration. The procedure is conservative (in simulation it
  accepted less than the best safe threshold), and its cost on dialogue comes from hand-coding whole transcripts and scaling
  every transcript by the largest one. Sampling single utterances avoids most of that, although at our budget the coders'
  confident dialogue decisions then still gave too little evidence. The population version additionally needs independent clusters and a true
  bound on cluster weight. For the corpus guarantee, random re-splits are a direct check (the test split is the uncoded
  remainder); for the population version they are only a diagnostic.
  \item \textbf{Post-hoc analysis.} The plan was not externally registered. Sonnet was extended to the full sample after subset
  results were seen, the A3 analyses were revised after their first output, and the corpus guarantee, the prevalence
  analysis and several controls were added or revised after external reviews. All are labelled exploratory.
  \item \textbf{LLM read-out.} Token probabilities were read after a generated rationale, and multi-label verbalized
  probabilities were converted under a reading the prompt did not fix (\S\ref{sec:setup}); both may disadvantage the LLM
  baselines. We tested one alternative prompt.
  \item \textbf{Adjudication} covered a stratified, high-confidence-enriched sample of 200 items drawn from disagreements
  with six label sources; it relied partly on an LLM panel whose families overlap with the coders it judged (Claude Opus with Haiku and
  Sonnet; GPT-6 Sol with Luna and OpenAI Decisions), and the author, the only human adjudicator, did not repeat a sitting to
  measure their own consistency.
  \item \textbf{Measurement.} RQ2 rests on four AnnoMI transcripts and is unweighted. ECE is top-label calibration against the
  majority label. Costs are October 2026 list prices for inference only; Kev's are estimated warm-serving costs.
  \item \textbf{Coverage gaps.} Kev-9B's robustness arms cover one to four tasks, and Kev-27B was run on the main and format
  arms only. Decision models cannot be probed for memorization; Kev's published training data excludes our datasets.
\end{itemize}

\section{Conclusion}
\label{sec:conclusion}
Model confidence can select low-error coding decisions, but less than accepted-decision shares suggest. Low calibration error
is not enough: what matters is how well confidence ranks errors and whether the top of the scale is resolved, and per-code
recalibration helps with both. Error targets need a guarantee, and a valid guarantee must respect how the hand-coded sample
was drawn. With one, models certified most decisions on Reddit comments, nearly all of them calls that a code does not
apply. On dialogue they certified nothing when whole transcripts were hand-coded; sampling single utterances made a
guarantee possible in principle, but at our budget the coders' confident dialogue decisions gave too little evidence for it. The best
decision-model configurations certified about as much as a frontier LLM on comments, at a small fraction of its cost.

\section*{Data and code availability}
The harness, the analysis plan with its amendments and deviations log, every model output (one row per decision, with the
probability and the model version returned), the samples and splits, the adjudication verdicts, a text-free spend ledger (405,339
requests; \$78.56 in API fees, plus GPU time for Kev) and the paper source are at \url{https://github.com/firish/margin}. The datasets are not redistributed; they are available from
their authors under their own terms (MFRC: CC BY 4.0; TalkMoves: CC BY-NC-SA 4.0; AnnoMI and Dunivin: no licence stated).

\section*{Ethics statement}
All data are public and were collected and annotated by their original authors; we used them only for evaluation. Only these
public texts were sent to commercial APIs; no private or newly collected data were used. AnnoMI's counselling conversations
are transcribed from publicly available demonstration videos of motivational interviewing. The only human participant was
the author, as adjudicator.

\section*{Use of AI assistants}
The author used an AI coding assistant (Claude) to implement the harness and to draft and edit the manuscript. An external
AI reviewer (GPT-6.1) commented on successive drafts, and its comments led to the changes logged as deviations 16--18. LLMs also served as
study subjects and as members of the adjudication panel, as described above. The author takes full responsibility for the
content.
